\subsubsection{Phase-Kickback Intuition}\label{sec:phase-kickback-intuition}

The remarkable feature of the Bernstein-Vazirani algorithm is that the \hl{oracle never explicitly writes the hidden bitstring $w$ into the data register.} Instead, it hides that information in the \textbf{phases} of the quantum state. This is possible because the ancilla is prepared in a very special state.

\highspace
\begin{flushleft}
    \textcolor{Green3}{\faIcon{flask} \textbf{The Ancilla is Prepared in \ket{-}}}
\end{flushleft}
After the first Hadamard gate, the ancilla is no longer in the computational basis:
\begin{center}
    \begin{quantikz}
        \lstick{$\ket{0}^{\otimes n}$}\slice[label style={inner sep=1pt,anchor=south west,rotate=40}]{\ket{\psi_0}} & & \qwbundle{n}  & \gate{H^{\otimes n}}\slice[label style={inner sep=1pt,anchor=south west,rotate=40}]{\ket{\psi_1}} & \\
        \lstick{$\ket{1}$}                                                                                          & &               & \gate{H}                                                                                                   &
    \end{quantikz}
\end{center}
Instead, it is in the state:
\begin{equation*}
    \ket{-} = \frac{\ket{0} - \ket{1}}{\sqrt{2}}
\end{equation*}
This state is an eigenstate of the Pauli-$X$ gate with eigenvalue $-1$:
\begin{equation*}
    X\ket{-} = -\ket{-}
\end{equation*}
Therefore, whenever a CNOT attempts to flip the ancilla, the ancilla itself is unchanged (i.e., it remains in the state $\ket{-}$), but the control state acquires a minus sign. This phenomenon is called \textbf{phase kickback} and is exactly the same mechanism encountered in Deutsch's algorithm (\autopageref{sec:deutsch-applying-the-oracle-to-the-superposition}). Let's see in detail how this works in the Bernstein-Vazirani algorithm.

\highspace
\textcolor{Green3}{\faIcon{cogs} \textbf{Oracle Action.}} The oracle always satisfies:
\begin{equation*}
    U_f\ket{x}\ket{y} = \ket{x}\ket{y \oplus f(x)}
\end{equation*}
But notice something important: the \hl{oracle \textbf{does not know} that the second register is in the state $\ket{-}$.} It simply receives \emph{some quantum state} as the second qubit and applies XOR. Since the Bernstein-Vazirani algorithm prepares the second register (ancilla) in the state \ket{-}, the input to the oracle is actually:
\begin{align*}
    \ket{x}\ket{-} = \frac{1}{\sqrt{2}}\left(\ket{x0} - \ket{x1}\right)
\end{align*}
Now, since the quantum gates are \textbf{linear}, we can rewrite the oracle action as:
\begin{align*}
    U_f\ket{x}\ket{-} &= U_f \left(\frac{1}{\sqrt{2}}\left(\ket{x0} - \ket{x1}\right)\right) \\
    &= \frac{1}{\sqrt{2}} \left(
        U_f\ket{x0} - U_f\ket{x1}
    \right) \\
    &= \frac{1}{\sqrt{2}} \left(
        U_f\ket{x}\ket{0} - U_f\ket{x}\ket{1}
    \right)
\end{align*}
At this point, we can apply the oracle definition to each term. Remember that, by definition, the oracle satisfies:
\begin{equation*}
    U_f\ket{x}\ket{y} = \ket{x}\ket{y \oplus f(x)}
\end{equation*}
Thus, if:
\begin{itemize}
    \item $y = 0$, the expression becomes:
    \begin{equation*}
        U_f \ket{x}\ket{0} = \ket{x} \ket{0 \oplus f(x)}
    \end{equation*}
    But XOR has a simple property:
    \begin{center}
        \begin{tabular}{@{} c c c @{}}
            \toprule
            $y$ & $f(x)$ & $y \oplus f(x)$ \\
            \midrule
            0 & 0 & 0 \\
            0 & 1 & 1 \\
            1 & 0 & 1 \\
            1 & 1 & 0 \\
            \bottomrule
        \end{tabular}
        $\quad y = 0 \implies y \oplus f(x) = f(x) \quad \forall f(x)$
    \end{center}
    Thus, we can rewrite the oracle action as:
    \begin{equation*}
        U_f \ket{x}\ket{0} = \ket{x} \ket{f(x)}
    \end{equation*}

    \item $y = 1$, the expression becomes:
    \begin{equation*}
        U_f \ket{x}\ket{1} = \ket{x} \ket{1 \oplus f(x)}
    \end{equation*}
    Similarly, the XOR property is the same, but on the opposite side:
    \begin{equation*}
        y = 1 \implies y \oplus f(x) = \overline{f(x)} \quad \forall f(x)
    \end{equation*}
    Thus, we can rewrite the oracle action as:
    \begin{equation*}
        U_f \ket{x}\ket{1} = \ket{x} \ket{\overline{f(x)}}
    \end{equation*}
\end{itemize}
Substituting these results back into the oracle action, we have:
\begin{align*}
    U_f\ket{x}\ket{-} &= \frac{1}{\sqrt{2}} \left(
        \ket{x}\ket{f(x)} - \ket{x}\ket{\overline{f(x)}}
    \right) \\
    &= \frac{1}{\sqrt{2}} \left(
        \ket{x}\ket{f(x)} - \ket{x}\ket{1 \oplus f(x)}
    \right)
\end{align*}
Leave $1 \oplus f(x)$ as it is for clarity. From this expression, we need to find a way to generalize it to all possible values (1 and 0) of $f(x)$:
\begin{itemize}
    \item $f(x) = 0$:
    \begin{align*}
        U_f\ket{x}\ket{-} &= \frac{1}{\sqrt{2}} \left(
            \ket{x}\ket{0} - \ket{x}\ket{1}
        \right) \\
        &= \frac{1}{\sqrt{2}} \left(
            \ket{x0} - \ket{x1}
        \right) \\
        &= \ket{x}\ket{-}
    \end{align*}
    \textbf{Nothing changes} in the state of the system, and the \hl{oracle does not introduce any phase change.}

    \item $f(x) = 1$:
    \begin{align*}
        U_f\ket{x}\ket{-} &= \frac{1}{\sqrt{2}} \left(
            \ket{x}\ket{1} - \ket{x}\ket{0}
        \right) \\
        &= -\frac{1}{\sqrt{2}} \left(
            \ket{x}\ket{0} - \ket{x}\ket{1}
        \right) \\
        &= -\frac{1}{\sqrt{2}} \left(
            \ket{x0} - \ket{x1}
        \right) \\
        &= -\ket{x}\ket{-}
    \end{align*}
    In this case, the \hl{oracle introduces a \textbf{phase change} of $-1$} in the state of the system.
\end{itemize}
Combining these two cases, we need to \hl{find a way to express the oracle action in a single expression that captures both cases:}
\begin{equation*}
    U_f\ket{x}\ket{-} = \begin{cases}
        \ket{x}\ket{-} & \text{if } f(x) = 0, \\
        -\ket{x}\ket{-} & \text{if } f(x) = 1.
    \end{cases}
\end{equation*}
The simplest way to express this is:
\begin{equation}\label{eq:bernstein-vazirani-oracle-action}
    U_f\ket{x}\ket{-} = (-1)^{f(x)}\ket{x}\ket{-}
\end{equation}
For the skeptics, we can verify:
\begin{gather*}
    f(x) = 0 \implies (-1)^{f(x)} = (-1)^0 = 1 \implies U_f\ket{x}\ket{-} = \ket{x}\ket{-} \\
    f(x) = 1 \implies (-1)^{f(x)} = (-1)^1 = -1 \implies U_f\ket{x}\ket{-} = -\ket{x}\ket{-}
\end{gather*}
Since we are very meticulous, we can expand $f(x)$ in the case of the Bernstein-Vazirani algorithm using the \autoref{eq:bernstein-vazirani-oracle-function} (\autopageref{eq:bernstein-vazirani-oracle-function}):
\begin{equation}\label{eq:bernstein-vazirani-oracle-action-expanded}
    U_f\ket{x}\ket{-} = \underbrace{(-1)^{w \cdot x}\ket{x}\ket{-}}_{\text{where } f(x) = w \cdot x \mod 2}
\end{equation}
Now, we can rewrite the Bernstein-Vazirani algorithm circuit (page \autopageref{fig:bernstein-vazirani-circuit}) with the explicit oracle action in mind:

\begin{figure}[!htp]
    \centering
    \begin{quantikz}
        \lstick{$\ket{0}^{\otimes n}$}\slice[label style={inner sep=1pt,anchor=south west,rotate=40}]{\ket{\psi_0}} & \qwbundle{n} & \gate{H^{\otimes n}}\slice[label style={inner sep=1pt,anchor=south west,rotate=40}]{\ket{\psi_1}} & \gate[2][5.21cm]{U_{f_w}}\gateinput{$\dfrac{1}{\sqrt{2}}\displaystyle\sum_{x} \ket{x}$}\gateoutput{$\dfrac{1}{\sqrt{2^{n}}}\displaystyle\sum_{x} \left(-1\right)^{f_w(x)}\ket{x} \ket{-}$}\slice[label style={inner sep=1pt,anchor=south west,rotate=40}]{\ket{\psi_2}}   & \qwbundle{n} & \gate{H^{\otimes n}}\slice[label style={inner sep=1pt,anchor=south west,rotate=40}]{\ket{\psi_3}} & \meter{} \\[1em]
        \lstick{$\ket{1}$}                                                                                          &              & \gate{H}                                                                                          & \gateinput{\ket{-}}\gateoutput{\ket{-}}                                                                                                                                                                                                                                &              &                                                                                                   & 
    \end{quantikz}
    \caption{Bernstein-Vazirani algorithm circuit with the explicit oracle action. The oracle is applied to the superposition of all possible inputs, and it introduces a phase change of $(-1)^{f_w(x)}$ to each basis state $\ket{x}$ in the superposition.}
    \label{fig:bernstein-vazirani-circuit-with-oracle-action}
\end{figure}

\newpage

\noindent
\begin{flushleft}
    \textcolor{Green3}{\faIcon{search} \textbf{Looking at One Bit at a Time}}
\end{flushleft}
The \autoref{eq:bernstein-vazirani-oracle-action-expanded} explains the entire algorithm mathematically, but it does not immediately reveal why the hidden bitstring $w$ eventually appears at the output. A much more intuitive way is to examine the oracle \textbf{bit by bit}.

\highspace
Recall that the oracle contains one CNOT gate for every position where $w_i = 1$ (\autopageref{sec:quantum-oracle-construction}), because the oracle computes the function:
\begin{equation*}
    f_w(x) = w \cdot x \mod 2 = w_0 x_0 \oplus w_1 x_1 \oplus \dots \oplus w_{n-1} x_{n-1}
\end{equation*}
Now consider \textbf{what happens to a single data qubit after the first\break Hadamard layer}. Every data qubit starts in:
\begin{equation*}
    \ket{+} = \frac{\ket{0} + \ket{1}}{\sqrt{2}}
\end{equation*}
Then, there are two possibilities: $w_i = 0$ or $w_i = 1$.

\highspace
\textcolor{Green3}{\faIcon{question-circle} \textbf{Why analyze the cases $w_i = 0$ and $w_i = 1$ when they are the goal and we cannot see them?}} Someone has already built the oracle. Suppose Nature secretly chooses a hidden bitstring:
\begin{equation*}
    w = 10110
\end{equation*}
The oracle is then physically built as:
\begin{center}
    \begin{quantikz}
        \lstick{$q_1$}          & \ctrl{5}  &             &           & \\
        \lstick{$q_2$}          &           &             &           & \\
        \lstick{$q_3$}          &           & \ctrl{3}    &           & \\
        \lstick{$q_4$}          &           &             & \ctrl{2}  & \\
        \lstick{$q_5$}          &           &             &           & \\
        \lstick{$\text{anc}$}   & \targ{}   & \targ{}     & \targ{}   & 
    \end{quantikz}
\end{center}
\noindent
Notice that the oracle contains a CNOT gate for every position where $w_i = 1$. The algorithm \textbf{never sees} this circuit, it simply receives the oracle as a black box.

\highspace
So we analyze the two cases $w_i = 0$ and $w_i = 1$ because \textbf{we}, as mathematicians, are \hl{trying to prove that the algorithm works \textbf{for every possible oracle}.}
\begin{itemize}
    \item Suppose the hidden bit at position $i$ is \hl{$w_i = 0$}. \emph{What happens?} There is no CNOT gate in the oracle for this position, so the \textbf{input data qubit $x_i$ remains unchanged}: \ket{+}.


    \item Now suppose instead that the hidden bit at position $i$ is \hl{$w_i = 1$}. \emph{What happens?} \textbf{There is a CNOT gate} in the oracle for this position.
    
    Let's suppose a \textbf{single data qubit} and one ancilla qubit.
    \begin{center}
        \begin{quantikz}
            \lstick{\ket{0}} & \gate{H} & \gate[2][1.7cm]{U_f}\gateinput{\ket{+}} \\
            \lstick{\ket{1}} & \gate{H} & \gateinput{\ket{-}}
        \end{quantikz}
    \end{center}
    Before the oracle, the state of the system is:
    \begin{align*}
        \ket{\psi_1} &= \ket{+}\ket{-} \\
        &= \frac{1}{\sqrt{2}}\left(\ket{0} + \ket{1}\right)\frac{1}{\sqrt{2}}\left(\ket{0} - \ket{1}\right) \\
        &= \frac{1}{2}\left(\ket{00} - \ket{01} + \ket{10} - \ket{11}\right)
    \end{align*}
    Now the oracle applies a CNOT gate:
    \begin{align*}
        \text{CNOT}\left(\ket{+}\ket{-}\right) &= \frac{1}{2}\left(\ket{00} - \ket{01} + \ket{11} - \ket{10}\right) \\
        &= \frac{1}{2}\left(\ket{00} - \ket{01} - \ket{10} + \ket{11}\right)
    \end{align*}
    Now factor the expression:
    \begin{align*}
        \text{CNOT}\left(\ket{+}\ket{-}\right) &= \frac{1}{2}\left(\ket{0}\left(\ket{0} - \ket{1}\right) - \ket{1}\left(\ket{0} - \ket{1}\right)\right) \\
        &= \frac{1}{2}\left(\ket{0} - \ket{1}\right)\left(\ket{0} - \ket{1}\right) \\
        &= \ket{-}\ket{-}
    \end{align*}
\end{itemize}
Therefore, \hl{each bit of the hidden bitstring $w$ determines whether the corresponding qubit ends up in the state $\ket{+}$ or $\ket{-}$ after the oracle.}
\begin{equation}
    x_i = \begin{cases}
        \ket{+} & \text{if } w_i = 0, \\
        \ket{-} & \text{if } w_i = 1.
    \end{cases}
\end{equation}
If $w_i = 0$, the qubit remains in $\ket{+}$; if $w_i = 1$, it flips to $\ket{-}$. This is the \textbf{essence of phase kickback}: the oracle encodes the hidden bitstring into the phases of the quantum state without ever explicitly revealing it in the data register.

\highspace
For example, if the hidden bitstring is $w = 10110$, corresponds to:
\begin{equation*}
    \ket{-} \otimes \ket{+} \otimes \ket{-} \otimes \ket{-} \otimes \ket{+}
\end{equation*}
The oracle has effectively transformed each hidden bit into a different quantum state of the corresponding data qubit.

\highspace
\begin{flushleft}
    \textcolor{Green3}{\faIcon{unlock-alt} \textbf{The Final Hadamard Decodes the Answer}}
\end{flushleft}
The last Hadamard layer \textbf{converts these phase states back into computational basis states}. Recall the identities:
\begin{align*}
    H\ket{+} &= \ket{0} \\
    H\ket{-} &= \ket{1}
\end{align*}
Thus, after the final Hadamard layer, the data qubits will be in the state
\begin{equation*}
    \ket{w_0}\ket{w_1}\ket{w_2}\ket{w_3}\ket{w_4} = \ket{10110}
\end{equation*}

\highspace
\begin{flushleft}
    \textcolor{Green3}{\faIcon{tools} \textbf{The Big Picture}}
\end{flushleft}
The Bernstein-Vazirani algorithm can now be summarized as a simple three-step process:
\begin{enumerate}
    \item The first Hadamard gates prepare every data qubit in the state \ket{+}.
    \item The oracle uses phase kickback to change selected qubits into the state \ket{-}.
    \item The final Hadamard gates convert:
    \begin{equation*}
        \ket{+} \to \ket{0} \quad \text{and} \quad \ket{-} \to \ket{1}
    \end{equation*}
    Revealing the hidden bitstring directly in the computational basis.
\end{enumerate}
This perspective explains why the algorithm succeeds with certainty: the oracle does not merely evaluate a Boolean function; it \textbf{encodes each hidden bit into the phase of its corresponding data qubit}, and the final Hadamard layer converts those phases back into ordinary bits.

\highspace
\begin{takeawaysbox}[: Phase Kickback Intuition]
    The ancilla is prepared in the state \ket{-}, so each CNOT in the oracle produces phase kickback instead of flipping the ancilla. As a result:
    \begin{equation*}
        \ket{x}\ket{-} \mapsto (-1)^{f_w(x) = w \cdot x}\ket{x}\ket{-}
    \end{equation*}
    Meaning that each hidden bit is encoded as a phase on the corresponding data qubit:
    \begin{equation*}
        w_i = 0 \implies x_i = \ket{+} \quad \text{and} \quad w_i = 1 \implies x_i = \ket{-}
    \end{equation*}
    The final Hadamard layer then decodes these phase states using:
    \begin{equation*}
        H\ket{+} = \ket{0} \quad \text{and} \quad H\ket{-} = \ket{1}
    \end{equation*}
    Deterministically recovering the hidden bitstring $w$ in the computational basis.
\end{takeawaysbox}